\documentclass[12pt,letterpaper]{article}
\usepackage[letterpaper,left=.88in,right=.88in,top=.83in,bottom=.80in,headheight=14pt]{geometry}
\usepackage{amsmath,amssymb,booktabs,tabularx,microtype,xcolor,fancyhdr}
\usepackage[colorlinks=true,urlcolor=teal!63!black,linkcolor=teal!63!black]{hyperref}
\definecolor{ink}{HTML}{163643}\definecolor{accent}{HTML}{116B66}
\pagestyle{fancy}\fancyhf{}
\fancyhead[L]{\small\sffamily\color{ink}INTEGER FACTORIZATION}
\fancyhead[R]{\small\sffamily\color{ink}RESEARCH NOTE\enspace/\enspace NOVEMBER 2026}
\fancyfoot[L]{\small\sffamily\color{ink}Divisor-sum bits of balanced semiprimes}
\fancyfoot[R]{\small\sffamily\color{ink}\thepage}
\renewcommand{\headrulewidth}{.4pt}
\setlength{\parindent}{0pt}\setlength{\parskip}{5pt}\setlength{\emergencystretch}{2em}
\newcommand{\sect}[1]{{\large\sffamily\bfseries\color{ink}#1}\par}
\newcommand{\subh}[1]{{\sffamily\bfseries\color{accent}#1}\par}
\newcommand{\F}{\mathbb F}\newcommand{\Z}{\mathbb Z}
\newcommand{\bits}{\operatorname{bits}}
\newcommand{\rs}[1]{\textup{[#1]}}
\newcommand{\src}[3]{\par\hangindent=1.8em\hangafter=1\noindent\makebox[1.8em][l]{\textsf{#1}}#2\enspace\href{#3}{\scriptsize\sffamily SOURCE}\par}
\hypersetup{pdftitle={Recovering divisor-sum bits of balanced semiprimes}}
\begin{document}
\smallskip
{\huge\bfseries\color{ink}Recovering divisor-sum bits\\[2pt]of balanced semiprimes\par}
\vspace{4pt}
{\large\sffamily\color{accent}Two exact recoveries and the limits of three classical searches\par}
\vspace{9pt}
A short block of bits of the divisor sum can identify both primes in a balanced semiprime. The difficult step is to obtain those bits from the integer. We recover two such blocks with no oracle input. A bounded elliptic-curve search succeeds on a 128-bit example. A weighted difference of squares succeeds on a 212-bit example missed by every curve in the same schedule. A separate 211-bit example defeats both searches. The inputs are synthetic and the two successful methods are classical \rs{1--3}.

\subh{Searches run on the integer alone}
\begin{tabularx}{\linewidth}{@{}l l l X@{}}
\toprule
\textbf{Input} & \textbf{Weighted squares} & \textbf{Elliptic curves} & \textbf{Recovered word}\\
\midrule
128 bits & Exhausted & Seed 19 & 48 bits, from index 0\\
212 bits & Multiplier 63 & Exhausted & 88 bits, from index 2\\
211 bits & Exhausted & Exhausted & No word recovered\\
\bottomrule
\end{tabularx}

Let $b=\bits(N)=\lfloor\log_2 N\rfloor+1$, let $\sigma$ sum the positive divisors, and set
\[
 B(z)=\left\lfloor 2^{\lfloor\bits(z)/2\rfloor}\frac{\sigma(z)}{z}\right\rfloor\bmod 2,
 \qquad z_i=2^{5b+2i}N .
\]
The earlier decoder chooses an index and a width from $N$, with $b^2$ candidate intervals. We use precisely those choices. Once a search splits $N=pq$, multiplicativity gives $\sigma(z_i)=(2^{5b+2i+1}-1)(p+1)(q+1)$. Reading $B(z_i)$ in increasing index order gives the following hexadecimal words, with four bits per digit.
\[
\begin{array}{rl}
128\text{-bit input}&\texttt{3587e5c0f45b}\\
212\text{-bit input}&\texttt{955b707abcaa6d41816c10}.
\end{array}
\]
The independent midpoint decoder returns the original factors in each case. The companion Lean development certifies the prime factors of the three inputs and the two word values for the original divisor-sum definition. The asymptotic results on later pages have mathematical proofs in this note.

\subh{What would settle the question}
A general factoring algorithm must find a proper divisor for every composite input within a fixed polynomial in $b$. The two schedules have polynomial bit cost. Their success depends on the input. Pages 4--5 identify an obstruction for a family of order tests on one fixed elliptic curve. Page 6 gives the exact cost in samples of a four-square collision method.
\newpage

\sect{1.\quad Two ways to obtain the bits}
\subh{Elliptic curves}
The first input is $N_1=199604481618592988980642015258549902581$. Use the Suyama parametrization in Montgomery coordinates \rs{1,2}. For each seed $s=6,\ldots,b+5$, set $u=s^2-5$, $v=4s$, $P=(u^3:v^3)$ and
\[
 A_{24}=(v-u)^3(3u+v)/(16u^3v) \pmod{N}.
\]
Test the denominator and $A_{24}(A_{24}-1)$ by gcd. Multiply $P$ by the largest power of each prime at most $b^2$, in increasing prime order, checking the projective $Z$ after each multiplication. The accumulated scalar is $\operatorname{lcm}(1,\ldots,b^2)$. For the second stage take $D=2\max(1,\lfloor\sqrt{b^3}/2\rfloor)$. For each prime $b^2<\ell\le b^3$, choose the nearest integer $j$ to $\ell/D$ (round ties up) and $r=|\ell-jD|$. Compare the $x$-coordinates of $[jD]Q$ and $[r]Q$ by cross multiplication, where $Q$ is the first-stage point.

For $N_1$, $s=19$ gives $D=1448$, $j=143$, $r=665$, and $\ell=206399$. The two projective $x$-coordinates agree modulo exactly one factor. The reduced cross product and its gcd give the following certificate.
\[
\begin{aligned}
 c&=175950992768636282823936990212460475462,\\
 \gcd(c,N_1)&=13326496357100595419,\\
 N_1/\gcd(c,N_1)&=14978016447079141999 .
\end{aligned}
\]
Multiplying $Q$ directly by $206399$ gives a $Z$-coordinate with the same proper gcd. Multiplication by $207729=jD+r$ gives a unit $Z$. All steps preceding the gcd use $N_1$ alone.

There are $b$ seeds and at most $b^3$ prime trials per seed. Each scalar explicitly passed to the Montgomery ladder has $O(\log b)$ bits, and each ladder step uses a constant number of operations modulo $N$. Sieving, gcds and screening therefore cost a polynomial in $b$ even when the entire schedule is exhausted.

\subh{Weighted differences of squares}
For $1\le k\le b^2$, start at $L_k=\lceil2\sqrt{kN}\rceil$ and test $A=L_k+j$ for every $0\le j\le\lfloor b^2/k\rfloor$. If $A^2-4kN=y^2$, test $\gcd((A-y)/2,N)$. The two numbers $A$ and $y$ have the same parity. This bounded form of the classical multiplier method uses $O(b^2\log b)$ square roots of $O(b)$-bit integers \rs{3}.

For the 212-bit input the first successful multiplier is $63$. Its first candidate satisfies
\[
 A=1022233444024044189366941464816376,\qquad A^2-4\cdot63N_2=698^2 .
\]
Gcd and division recover
\[
\begin{aligned}
 p_2&=56790746890224677187052303600871,\\
 q_2&=73016674573146013526210104629791 .
\end{aligned}
\]
Here $7q_2-9p_2=698$. The search received neither 63 nor the ratio. The test input used primes near $0.7\cdot 2^{106}$ and $0.9\cdot 2^{106}$, which explains its favorable ratio.
\newpage

\sect{2.\quad When a small multiplier or a point count suffices}
\subh{Range of the multiplier method}
Suppose $N=pq$ and $k=uv$. The relevant difference is
\[
 (up+vq)-2\sqrt{kN}=\frac{(up-vq)^2}{up+vq+2\sqrt{kN}} .
\]
If $1\le k\le b^2$ and $0\le up+vq-L_k\le\lfloor b^2/k\rfloor$, the search tests $A=up+vq$. If $p,q>k$, the half-difference $(A-|up-vq|)/2$ yields a proper gcd.

Conversely, consider primes $p,q\in[X,2X]$, let $L=\lceil\log_2(2X)\rceil$, and allow any multipliers and offsets bounded by $L^C$ for fixed $C>0$. Every success produces two positive integers $(A-y)/2,(A+y)/2$ with product $kN$. For sufficiently large $X$, $p,q>k$. A proper gcd forces the two integers to be $up,vq$ in some order with $uv=k$. Since $A\le2\sqrt{kN}+1+L^C$,
\[
 |up-vq|^2=A^2-4kN=O_C(XL^{2C+2}).
\]
For each $p$ and $u,v\le L^C$ this leaves $O_C(\sqrt X\,L^{C+1})$ integer values of $q$. Summing over $u,v$ and dividing by the $\asymp X/L$ primes in the interval gives a success fraction $O_C(X^{-1/2}L^{3C+2})$. This bound concerns precisely square tests within those ranges. Lehman's full algorithm uses larger multipliers.

\subh{An input missed by both searches}
For the independently generated 211-bit semiprime
\[
 N_3=2326981664014727582718789409779272096176399345327619480698638841
\]
every one of the 527,954 square tests and all 211 curve seeds completed without a proper gcd. Its prime factors are $42678434690689936321894022265043$ and $54523594430756986474678777350787$. Their primality and product certify that $N_3$ meets the input promise. The computation refutes completeness for the stated combination of schedules.

\subh{A fixed curve whose order is the desired divisor sum}
Suppose now that $p,q$ are distinct primes congruent to $3\bmod4$ and $E$ is $y^2=x^3-x$. For $\ell\in\{p,q\}$ write $\chi$ for the quadratic character modulo $\ell$, with $\chi(0)=0$. The summands for $x$ and $-x$ in $\sum_x\chi(x^3-x)$ cancel because $\chi(-1)=-1$. Thus $E$ has exactly $\ell$ affine points and $\ell+1$ projective points. The Chinese remainder theorem gives
\[
 \#E_{\rm aff}(\Z/N\Z)=N,\qquad
 \#E(\Z/N\Z)=(p+1)(q+1)=\sigma(N).
\]
The projective point set is $E(\F_p)\times E(\F_q)$. Outside the affine chart are the common point at infinity and $p+q$ points infinite in just one component. For any such mixed point, a primitive projective representative $(X:Y:Z)$ satisfies $1<\gcd(Z,N)<N$. An exact projective count $C$ also gives $s=C-N-1=p+q$ and hence the factors $(s\pm\sqrt{s^2-4N})/2$. The trace-zero affine count is recorded explicitly in the point-count literature \rs{4}. The mixed points account for all the information missing from that count.
\newpage
\sect{3.\quad A restriction on order tests for the fixed curve}
Keep $E:y^2=x^3-x$. For any $x\in\F_\ell$ and $\ell\equiv3\bmod4$ prime, take a point $P_\ell$ with that $x$-coordinate on $E$ over an algebraic closure. The two choices of sign have the same order. This order divides $\ell+1$. Indeed, either the point is rational or it corresponds to a rational point of the nonsquare quadratic twist. The twist has the same trace zero by the character sum on page 3. This argument allows arbitrary starting $x$-coordinates.

\subh{Theorem}
Fix $k>0$. At each scale $X$, fix two lists of integer polynomials $h$ and $f$, each of length at most $L^k$, where $L=\lceil\log_2(2X)\rceil$. Each polynomial has degree at most $L^k$ and coefficients of bit length at most $L^k$. The lists may depend on $X$. Write $T=\lfloor\sqrt{pq}\rfloor$. Choose a start $x=h(T)$ and any positive exponent $e$ whose prime divisors are at most $L^k$ or divide some nonzero value $f(pq)$ or $f(T)$ from the list. The choices of $h$ and $e$ may depend on the input.

For every fixed $s>0$, apart from an $O_{k,s}(L^{-s})$ fraction of ordered pairs of distinct primes $p,q\in[X,2X]$, both $3\bmod4$, the following equivalence holds simultaneously for every permitted $h,e$:
\[
        [e]P_p=O \quad\Longleftrightarrow\quad [e]P_q=O .
\]
A test that splits by killing exactly one component has negligible coverage in this model. The exponent may be arbitrarily large. Taking $f(t)=t-1$ or $t+1$ permits all prime factors of $N\pm1$. Input-dependent curves with varying $j$-invariant, unrestricted starting-point circuits, and other ways of extracting gcds fall outside the statement.

The theorem also controls two-stage comparisons on this fixed curve. If $a>b>0$, equality of the projective $x$-coordinates of $[aM]P$ and $[bM]P$ means that $[M(a-b)]P=O$ or $[M(a+b)]P=O$. When $M$ has allowed prime support and $a+b$ is bounded by a fixed power of $L$, both tests fit the theorem after increasing $k$ if necessary. The run in section 1 varies the curve itself. Its 211-bit counterexample was obtained by direct computation.

\subh{Proof, part I. Prime factors allowed in the exponent}
Put $K=2k+s+8$, $B=L^K$ and $Y=X^{1/10}$. The $B$-smooth part of an integer is the product of all its prime powers with base at most $B$. For almost all eligible $p$, the $B$-smooth part of $p+1$ is less than $Y$. Indeed, if it reaches $Y$, multiply the individual prime factors, with multiplicity, up to the first crossing. The resulting divisor $w$ lies in $[Y,BY]$. For each $w$, at most $X/w+1\le2X/w$ integers in the interval are $-1\bmod w$, for large enough $X$. For $\delta=1/(2(K+1))$,
\[
 \sum_{\substack{w\ge Y\\w\ B\text{-smooth}}}\frac1w
 \le Y^{-\delta}\prod_{\substack{\ell\le B\\\ell\ {\rm prime}}}(1-\ell^{-1+\delta})^{-1}
 \le Y^{-\delta}\exp(O_\delta(B^\delta)).
\]
The last estimate follows by summing $O_\delta(m^{-1+\delta})$ up to $B$. Since $B^\delta=o(\log X)$, the exceptional integer count is $X^{1-\delta/10+o(1)}$. The prime number theorem in the progression $3\bmod4$ gives $\asymp X/L$ eligible primes. This error is smaller than every fixed inverse power of $L$ as a fraction of those primes. Smooth parts of shifted primes are also studied in \rs{5}.
\newpage
\sect{4.\quad Finishing the order-test proof}
\subh{Large primes from polynomial values}
Discard identically zero $f$ and omit zero evaluations. The contents of all remaining $f$ have $O(L^{2k})$ distinct prime divisors in total, by their bit lengths. For each such prime $\ell>B$, discard $p\equiv-1\bmod\ell$. This costs at most $O(L^{2k}(X/B+1))$ integers $p$. For each remaining $p$, there are $O(L)$ distinct prime divisors of $p+1$. At any such $\ell>B$, every listed $f$ reduces to a nonzero polynomial. Since $p$ is a unit modulo $\ell$, the condition $\ell\mid f(pq)$ restricts $q$ to at most $\deg f$ residues. For $\ell\mid f(T)$ use the same count for $T\in[X,2X]$. For fixed $p,T$, the inequalities $T^2\le pq<(T+1)^2$ allow at most five integer $q$ for large $X$. At most $O(L^{2k+1}(X/B+1))$ values of $q$ are lost per $p$. Remove the symmetric exceptional pairs. Dividing by the prime counts gives an exceptional fraction $O_{k,s}(L^{-s})$.

On the surviving pairs, an order dividing $e$ and $p+1$ is $B$-smooth and less than $Y$. It remains to transfer any such short order from $p$ to $q$.

\subh{Short orders transfer between the primes}
For $1\le r<Y$ let $D_r(U)\in\Z[U]$ vanish exactly on the finite $x$-coordinates of $r$-torsion in odd characteristic $\ell>r$. To construct it, write the standard division polynomial $\psi_r$ \rs{6} as $F_r$ for odd $r$ and $2yF_r$ for even $r$. Take $D_r=F_r$ for odd $r$ and $D_r=(U^3-U)F_r$ for even $r$, including the two-torsion. Both degree and $\log(1+\|D_r\|_1)$ are $O(r^4)$, where the norm sums absolute coefficients. In fact the recurrences at index $r$ use differences of products of at most four $F_i$ with $i\le\lfloor r/2\rfloor+2$. For odd indices one product may have the extra factor $16(U^3-U)^2$. For even indices the normalized recurrence is
\[
 F_{2m}=F_m(F_{m+2}F_{m-1}^2-F_{m-2}F_{m+1}^2).
\]
Each size measure up to $r$ is therefore at most four times the maximum at $\lfloor r/2\rfloor+2$, plus a constant. For $r\ge16$ this index is at most $5r/8$. Since $4(5/8)^4<1$, induction gives the stated bound.

For each $h,r$ such that $D_r(h(t))$ is nonzero as an integer polynomial, select one nonzero coefficient. Across these polynomials the sum of degrees and of selected coefficient bit lengths is $O(Y^5 L^{O_k(1)})=X^{1/2+o(1)}$ by composition and the coefficient norm bound. Exclude primes dividing any selected coefficient. For remaining $p$, each polynomial has at most its degree many roots modulo $p$. Each root accounts for at most three possible $T$ in the interval and at most five $q$ per $T$. Discard these pairs in both orders. Their relative proportion is $X^{-1/2+o(1)}$.

If $[e]P_p=O$, its order $r$ divides $e$ and is less than $Y$. Now $D_r(h(T))=0\bmod p$ forces $D_r(h(t))$ to vanish identically by the exclusions. It vanishes modulo $q$ at $T$, giving $[r]P_q=O$ and hence $[e]P_q=O$. Interchanging $p,q$ proves the equivalence. \hfill$\square$
\newpage
\sect{5.\quad Four squares and independent samples}
For distinct odd primes $p,q$, Jacobi's identity gives exactly $8(p+1)(q+1)$ signed, ordered four-square representations of $N=pq$ \rs{7}. Algorithms for producing one representation are classical \rs{8}. Suppose a perfect independent uniform sampler is available. Interpret each representation as an integral Hamilton quaternion. Compare two samples $\alpha,\beta$ by taking the gcd of $N$ and the four coordinates of $\beta\bar\alpha$.

\subh{Exact success probability}
For two independent samples,
\[
 \Pr(\text{proper gcd})=\frac{p+q}{(p+1)(q+1)} .
\]
The Hamilton algebra modulo an odd prime $\ell$ identifies with $M_2(\F_\ell)$, with norm equal to determinant and conjugation equal to adjugation. Explicitly take
$i=\left(\begin{smallmatrix}0&1\\-1&0\end{smallmatrix}\right)$ and
$j=\left(\begin{smallmatrix}u&v\\v&-u\end{smallmatrix}\right)$,
where $u^2+v^2=-1$. Such $u,v$ exist because the squares and minus one minus the squares each occupy $(\ell+1)/2$ residues. The four matrices $1,i,j,ij$ form a basis and satisfy the quaternion relations. Related projective-line descriptions appear in \rs{9}.

Each representation reduces to rank one at $p$ and $q$ since divisibility of all coordinates by $\ell$ would give $\ell^2\mid N$. At either prime, $\beta\bar\alpha=0$ precisely when the one-dimensional right kernels of $\alpha$ and $\beta$ coincide. If both kernels coincide, $\beta\bar\alpha/N$ is integral with norm one. Left multiplication by it takes $\alpha$ to $\beta$. Conversely left multiplication by any of the eight integral units preserves both kernels and acts freely. Thus unit orbits inject into $\mathbb P^1(\F_p)\times\mathbb P^1(\F_q)$. Jacobi's count gives exactly $(p+1)(q+1)$ orbits, so the injection is a bijection. A proper gcd means agreeing in exactly one coordinate. Counting the two possibilities gives the formula.

\subh{Number of samples required}
For $m$ independent samples and all pairwise tests, the union bound gives success probability at most $\binom m2(p+q)/((p+1)(q+1))$. A first-coordinate collision occurs with probability at least $1-\exp(-m(m-1)/(2(p+1)))$ for $m\le p+1$. Subtract at most $\binom m2/((p+1)(q+1))$ for any pair agreeing in both coordinates. These bounds give $\Theta(N^{1/4})$ samples for constant success probability on balanced inputs.

An exact no-success probability is
\[
 \frac{\displaystyle\sum_{t=0}^{\min(m,p+1,q+1)}
 {m\brace t}(p+1)_t(q+1)_t}{((p+1)(q+1))^m}.
\]
Here \({m\brace t}\) counts set partitions and $(a)_t$ is a falling factorial. With no proper gcd, the partitions induced by equal first and equal second coordinates must coincide. A partition with $t$ blocks has $(p+1)_t(q+1)_t$ coordinate assignments. The argument leaves open what a construction of correlated representations could achieve.
\newpage

\sect{6.\quad Evidence, scope and references}
\subh{Reproduction}
The companion includes the two integer-only searches, the earlier decoder, the Lean project, and exact per-input results. \texttt{python verify.py} checks the numerical certificates and small exhaustive four-square tests. \texttt{python verify.py --replay} reruns both schedules on all three inputs. The searches accept only $N$. The known factors serve to construct or verify fixtures. The third fixture uses SHA-256 labels and the next primes congruent to $3\bmod4$. For each input, both factors are prime and the larger is less than twice the smaller.

\subh{Relation to recent preprints}
Shunia \rs{10} reduces general factoring to a polynomial number of abundance decisions. The cost of answering those decisions remains part of an unconditional algorithm. Gao and Feng \rs{11} use a second lattice vector to improve the logarithmic factors in a balanced-semiprime bound that still contains $N^{1/5}$. A November 2026 Cogentic manuscript \rs{12} reports deterministic general splitting in $O(\exp(3(\ln N)^{3/4}))$ bit operations. In binary input length this reported bound grows faster than every polynomial. Our examples give finite coverage results and mathematical limitations for the specified methods. The general worst-case polynomial bound remains open.

\subh{References}
{\small\setlength{\parskip}{3pt}
\src{1}{H. W. Lenstra, Jr. \emph{Factoring integers with elliptic curves}. Annals of Mathematics, 1987.}{https://annals.math.princeton.edu/1987/126-3/p09}
\src{2}{R. P. Brent. \emph{Factorization of the tenth Fermat number}. Mathematics of Computation, 1999, sections 2--3.}{https://maths-people.anu.edu.au/brent/pd/rpb161.pdf}
\src{3}{R. S. Lehman. \emph{Factoring large integers}. Mathematics of Computation, 1974.}{https://doi.org/10.1090/S0025-5718-1974-0340163-2}
\src{4}{L. V. Dieulefait and J. Urroz. \emph{Factorization and malleability of RSA moduli, and counting points on elliptic curves modulo N}. Mathematics, 2020, Remark 2.}{https://doi.org/10.3390/math8122126}
\src{5}{C. Pomerance and I. E. Shparlinski. \emph{Smooth orders and cryptographic applications}. ANTS V, 2002.}{https://math.dartmouth.edu/~carlp/PDF/paper128.pdf}
\src{6}{A. V. Sutherland. \emph{Elliptic curves}. MIT 18.783, division polynomials and multiplication by $n$.}{https://ocw.mit.edu/courses/18-783-elliptic-curves-spring-2021/59e0623c0d3454b39d8e07b493a0e9b0_MIT18_783S21_Slides5.pdf}
\src{7}{E. Assing. \emph{Advanced Algebra I: Modular Forms}, Theorem 4.0.3, Jacobi's four-square identity.}{https://www.math.uni-bonn.de/people/assing/lectures/modular_forms.pdf}
\src{8}{P. Pollack and E. Trevi\~no. \emph{Finding the four squares in Lagrange's theorem}. Integers, 2018.}{https://pollack.uga.edu/finding4squares.pdf}
\src{9}{A. Forsyth, J. Gurev and S. Shrima. \emph{Metacommutation as a group action on the projective line over $\F_p$}. Preprint, 2015.}{https://arxiv.org/abs/1503.06259}
\src{10}{J. M. Shunia. \emph{Deterministic integer factorization from abundance decisions}. Preprint, 2026.}{https://arxiv.org/abs/2611.34242}
\src{11}{Y. Gao and Y. Feng. \emph{On factoring and power divisor problems via rank-3 lattices and the second vector}. Preprint, 2025.}{https://arxiv.org/abs/2512.19076}
\src{12}{Cogentic. \emph{Unconditional deterministic subexponential integer factoring and quadratic nonresidue construction}. Research manuscript, 2026.}{https://github.com/google-research/google-research/tree/master/cogentic_research_papers/deterministic_factoring_l34}
}
\end{document}

